\documentclass[10pt,twocolumn]{article}
\usepackage[T1]{fontenc}
\usepackage[utf8]{inputenc}
\usepackage{lmodern}
\usepackage[margin=1.8cm,top=2.4cm,bottom=2cm,columnsep=0.6cm]{geometry}
\usepackage{fancyhdr}
\usepackage{titlesec}
\usepackage{booktabs}
\usepackage{amsmath,amssymb,amsfonts}
\usepackage{microtype}
\usepackage{xcolor}
\usepackage{mdframed}
\usepackage{parskip}
\usepackage{enumitem}
\usepackage{hyperref}

\definecolor{rulered}{RGB}{180,30,30}
\definecolor{boxgray}{RGB}{245,245,245}
\definecolor{keywordgray}{RGB}{100,100,100}

\pagestyle{fancy}
\fancyhf{}
\fancyhead[L]{\textsc{\small Claude Mag}}
\fancyhead[C]{\small\textit{Machine Learning Research Digest}}
\fancyhead[R]{\small 2026-07-20}
\fancyfoot[C]{\small\thepage}
\renewcommand{\headrulewidth}{0.8pt}

\titleformat{\section}{\normalsize\bfseries\color{rulered}}{}{0pt}{}%
  [\vspace{-4pt}\color{rulered}\rule{\columnwidth}{0.4pt}]
\titlespacing{\section}{0pt}{8pt}{4pt}

\hypersetup{colorlinks=true,urlcolor=rulered,linkcolor=black}

\begin{document}
\setlength{\parindent}{0pt}
\setlength{\parskip}{4pt}

{\small\color{keywordgray}\textbf{Keywords:}
  unified vision \textbullet{}
  RGB encoding \textbullet{}
  dense perception \textbullet{}
  zero-shot transfer}\par\vspace{4pt}

{\LARGE\bfseries\raggedright
  RINO}\par\vspace{4pt}

{\large\itshape\raggedright
  A Multi-Institution Team Proposes Treating All Vision Tasks
  as RGB-to-RGB Image Editing, Enabling One Backbone to Handle
  Depth, Segmentation, and Normals Without Fine-Tuning}\par\vspace{6pt}

{\small\textbf{By} Timing Yang, Jinrui Yang, Xinlong Li, Yuhan Wang,
  Haoran Li, Yanqing Liu, Guoyizhe Wei, Jixuan Ying, Chen Wei,
  Rama Chellappa, Yuyin Zhou, Cihang Xie, Alan Yuille, Feng Wang
  (Johns Hopkins University, UC Santa Cruz, CMU, Rice
  University)}\par\vspace{2pt}

{\small\color{keywordgray}arXiv:2607.12450 \textbullet{}
  July 14, 2026}
\par\vspace{2pt}\noindent{\color{rulered}\rule{\columnwidth}{0.6pt}}\vspace{6pt}

The dominant paradigm in computer vision maintains separate models for
separate tasks: one model for depth, another for segmentation, a third
for surface normals. Even modern unified models typically require
task-specific heads or adapters and must be fine-tuned per benchmark.
RINO (RGB In, RGB Out) proposes a different architecture: encode every
visual output as an RGB image and reduce every vision task to an
RGB-to-RGB image editing problem, sharing all parameters across all
tasks without task-specific components.

\begin{mdframed}[backgroundcolor=boxgray,linecolor=rulered,linewidth=1pt,
    innerleftmargin=6pt,innerrightmargin=6pt,
    innertopmargin=6pt,innerbottommargin=6pt]
\textbf{\small AT A GLANCE}\par\vspace{4pt}
\begin{description}[leftmargin=0pt,labelsep=4pt,itemsep=2pt,parsep=0pt,
    font=\small\bfseries]
  \item[Problem:] Dense visual perception requires separate models per
    task; existing unified approaches still use task-specific heads or
    require benchmark fine-tuning.
  \item[Why it matters:] True generalist vision requires a single set
    of parameters that transfers across tasks through a unified
    interface, as text does for NLP.
  \item[Method:] All visual modalities (depth, masks, normals,
    keypoints) encoded as RGB images; tasks become text-conditioned
    RGB-to-RGB editing problems.
  \item[Result:] Competitive zero-shot performance on four dense
    perception benchmarks using a single generic image editing backbone.
  \item[Evidence:] NYU-Depth-v2, ADE20K segmentation, ScanNet normals,
    H3WB keypoints; zero-shot, no benchmark fine-tuning.
\end{description}
\end{mdframed}

\section{The Big Picture}

Language models achieve generality through a single universal
representation: text. Any NLP task --- translation, summarization,
QA, classification --- can be expressed as a text-in, text-out
problem, allowing one set of parameters to transfer across all of
them. RINO asks whether pixels can serve the same function for
vision.

The central observation is that all dense vision outputs are
already spatial arrays over the same 2D grid as the input image.
Depth is a scalar per pixel, surface normals are a 3-vector per
pixel, segmentation masks are a class label per pixel, and keypoint
heatmaps are a confidence per pixel. All of these can be rendered
as an RGB image and decoded back without loss, if the encoding
is chosen carefully.

\section{Why Existing Approaches Fall Short}

Task-specific models are powerful but expensive: each benchmark
requires a separate model, separate fine-tuning run, and separate
inference process. Multi-task models (e.g., UniPerceiver, InternImage)
unify multiple tasks in one architecture but still require task-specific
output heads that must be trained on benchmark data. Generalist models
like Painter and PromptDiffusion reduce fine-tuning requirements but
retain task-specific structured output formats.

The key gap is that no prior work treats the task output itself as
RGB data fed into a backbone that was not designed for the task at
all. RINO demonstrates that a backbone trained entirely on natural
image editing --- with no exposure to depth maps, segmentation masks,
or surface normals --- can perform these tasks zero-shot if the output
is presented to it as an RGB image to edit.

\section{Core Mechanism}

Each visual modality is assigned a lossless or near-lossless RGB
encoding:

\begin{enumerate}[itemsep=2pt,parsep=0pt]
  \item \textbf{Depth maps:} Depth values $d$ are normalized to
    $[0, 255]$ and stored in the red channel. The green channel
    encodes epistemic uncertainty; blue is zero.
  \item \textbf{Segmentation masks:} Each semantic class receives a
    fixed, perceptually distinct RGB color. Instance masks use
    hue-shifted variants of the class color, enabling visual
    discrimination of instances.
  \item \textbf{Surface normals:} The XYZ components of the unit
    normal vector are linearly mapped to $[0, 255]$ triplets:
    $\text{RGB} = \lfloor 127.5(\mathbf{n} + \mathbf{1}) \rfloor$.
  \item \textbf{Keypoint heatmaps:} Per-joint Gaussian confidence
    heatmaps are stacked into the RGB channels (up to three joints
    per RGB image, with multiple images for larger joint sets).
\end{enumerate}

A text instruction specifies the desired output type. The editing
backbone processes the input image and instruction, producing an RGB
output that is decoded into the task-specific representation.

\section{Technical Formulation}

For a task $T$ with input image $\mathbf{x} \in \mathbb{R}^{H \times W \times 3}$
and ground truth output $\mathbf{y}_T$, RINO defines an encoding
function $\phi_T : \mathbf{y}_T \to \mathbb{R}^{H \times W \times 3}$
and decoding function $\phi_T^{-1}$. The backbone $f_\theta$ is trained
to minimize:
\[
\mathcal{L} = \mathbb{E}_{(x, y_T, T)} \left\|
  f_\theta(\mathbf{x},\, \tau_T) - \phi_T(\mathbf{y}_T)
\right\|_2^2
\]
where $\tau_T$ is the text instruction for task $T$.
Crucially, $f_\theta$ is the pretrained image editing backbone ---
its weights are frozen for zero-shot evaluation.

\section{Experimental Findings}

\begin{center}
\small
\begin{tabular}{llcc}
\toprule
Task & Metric & RINO & Fine-Tuned \\
\midrule
Depth (NYU-v2) & AbsRel$\downarrow$ & 0.137 & 0.103 \\
Seg. (ADE20K) & mIoU$\uparrow$ & 38.7 & 47.2 \\
Normals (ScanNet) & Mean Ang.$\downarrow$ & 11.3\textdegree & 9.1\textdegree \\
Keypoints (H3WB) & PCK$\uparrow$ & 61.4 & 68.9 \\
\bottomrule
\end{tabular}
\end{center}

RINO's zero-shot numbers are within 15--25\% (relative) of dedicated
fine-tuned models on each task, using a single set of backbone
weights that was never trained on any of these datasets.

\section{Ablations and Interpretation}

\textbf{Task-specific heads:} Adding a lightweight task-specific
head (while keeping the backbone frozen) closes approximately half
the gap to fully fine-tuned models on depth and segmentation.
This confirms that the backbone captures the necessary spatial
representations; the remaining gap is in task-specific prediction.

\textbf{RGB encoding quality:} Using a uniform grayscale depth
encoding (instead of the normalized red-channel scheme) degrades
depth AbsRel by 0.021 absolute, confirming that encoding choices
matter substantially.

\textbf{Instruction specificity:} Vague instructions (``transform
this image'') yield near-random outputs; specific instructions
(``generate a surface normal map with outward normals colored
red-green-blue'') approach zero-shot quality, suggesting the backbone
uses the instruction as a strong conditioning signal.

\textbf{Joint training benefit:} When the backbone is fine-tuned
jointly across depth + normals + segmentation (vs.\ separately),
each task improves by 1--2\%, suggesting mutual regularization from
shared RGB representations.

\section*{Reference}

Timing Yang et al.\
\textbf{Let RGB Be the Language of Vision}.
arXiv:2607.12450, July 2026.
\url{https://arxiv.org/abs/2607.12450}

\end{document}
